\subsection{THETA3-C22: the weaker slope is sharp for an actual tame theta section}\label{C22}

The stronger integral-function estimate in C19 cannot be extended
with a bounded loss to every individual tame reduced theta section.
The following explicit example determines the optimal possible
linear slope for that proposed extension.

\paragraph{Statement with the actual theta normalization.}
Take
\[
 E:\quad y^2=4x^3-4x,\qquad p=7,\qquad
 \omega=dx/y,\qquad t=-2x/y,
 \qquad T=(0,0).
\]
This curve has good reduction at 7 and CM by $\Z[i]$, and 7 is
inert in $\Q(i)$.  The point $T$ is a nonzero point of order
two.  Let $z_T$ be a half-period representing $T$, with a
compatible torsion lift as in C18, and use the actual Fock
translation
\[
 \sigma_{z_T}(z)=
 e^{-\eta_{\rm per}(z_T)(z+z_T/2)}\sigma(z+z_T).
\]
Here the real-linear quasiperiod map is extended to the
half-period in the indicated formula.  The period lattice is
$i$-stable, so $e_2^*=0$ and this is already the reduced theta
translation.  Normalize its nonzero constant by putting
\[
 F(z)=\frac{\sigma_{z_T}(z)}{\sigma_{z_T}(0)},\qquad
 f_T(t)=F(\lambda(t)),\qquad F(0)=1.
\]
The algebraic value $\sigma_{z_T}(0)$ is a 7-adic unit by
Katagiri--Kobayashi, arXiv:2609.07024v1,
Theorem~2.7, pp.~7--8, and Lemma~3.3(1), pp.~9--10.
Thus this normalization does not change any derivative valuation.
If $D$ is the invariant derivation dual to $\omega$, then for
every integer $r\geq0$, with $n_r=8\cdot7^r$,
\begin{equation}\label{C22-exact}
 v_7\bigl(D^{n_r}f_T(0)\bigr)
      =\frac{7^r-1}{6}
      =\frac{n_r}{48}-\frac16.
\end{equation}
Consequently no bound
$v_7(D^nf_T(0))\geq c n-C$ with $C$ independent of $n$
can hold for any $c>1/48$.
In particular the coefficient $7/48$ from the degree-zero
integral-function theorem cannot hold for this degree-one
theta section, even with an arbitrarily large fixed loss.

\paragraph{Evenness and the logarithmic derivatives.}
The sigma normalization and quasiperiod identities are those of
Bannai--Kobayashi, arXiv:math/0610163v4,
Example~1.9, pp.~10--11.
The automorphism $(x,y)\mapsto(-x,iy)$ pulls $\omega$ back to
$i\omega$, so the lattice is $i$-stable.  The weight-two
transformation of $e_2^*$, or comparison in
$\eta_{\rm per}(\gamma)=e_2^*\gamma+C\bar\gamma$, gives
$e_2^*=-e_2^*$ and hence $e_2^*=0$.

Since $2z_T$ is a lattice element not divisible by two in the
lattice, the sigma multiplier at $2z_T$ is $-1$.
The translation formula and oddness of sigma therefore give
\[
 \sigma(z+z_T)
     =-e^{2\eta_{\rm per}(z_T)z}\sigma(z-z_T),
 \qquad
 \sigma(z_T-z)
     =e^{-2\eta_{\rm per}(z_T)z}\sigma(z+z_T).
\]
After multiplication by the defining exponential factor these
identities say $F(-z)=F(z)$.  Thus $\log F$ is even, with
constant and linear terms zero.
The second derivative has the particularly simple algebraic form
\begin{equation}\label{C22-B}
 \frac{d^2}{dz^2}\log F(z)=-\wp(z+z_T).
\end{equation}
The right side is $-x(T+P(t))$ in the formal parameter at the
identity.  It belongs to $R[[t]]$ after a finite extension
carrying the torsion data: translation by $T$ is integral,
and its residue disc is disjoint from the sole pole of $x$,
because $T$ is nonzero prime-to-7 torsion.

Put $l_j=(d/dz)^j\log F(0)$.  Equation \eqref{C22-B}
and the intrinsic congruence $D^7R[[t]]\subset7R[[t]]$ of
C19 give
\begin{equation}\label{C22-log-bound}
 l_j=0\quad(j\text{ odd}),\qquad
 v_7(l_j)\geq\left\lfloor\frac{j-2}{7}\right\rfloor
       \quad(j\geq2\text{ even}).
\end{equation}
The decisive logarithmic derivative is a unit.  Indeed,
$Dx=y$ and $Dy=6x^2-2$ imply by direct differentiation
\[
 D^2x=6x^2-2,\qquad
 D^4x=120x^3-72x,\qquad
 D^6x=5040x^4-4032x^2+144.
\]
Evaluating at $T$ gives
\begin{equation}\label{C22-unit}
 l_8=-D^6x(T)=-144\in R^\times.
\end{equation}
The same computation gives $l_2=0$ and $l_4=2$, although
only \eqref{C22-unit} is needed for sharpness.

\paragraph{The unique lowest-valuation Bell term.}
Set $s=7/48$.  For a nonzero even logarithmic derivative
define the surplus
\[
 \epsilon_j=v_7(l_j)-v_7(j!)+js-\frac16.
\]
We have $\epsilon_8=0$, and
\begin{equation}\label{C22-strict}
 \epsilon_j>0\qquad(j\ne8\text{ even},\ l_j\ne0).
\end{equation}
Here is the strict inequality in full, rather than an appeal
only to the non-strict Bell bound of C18.
For $j=2,4,6$, \eqref{C22-log-bound} gives
$\epsilon_j\geq(7j-8)/48>0$.
For every even $j\geq10$, use
\[
 v_7(j!)\leq\frac{j-1}{6},\qquad
 \left\lfloor\frac{j-2}{7}\right\rfloor
       \geq\frac{j-8}{7}.
\]
It follows that
\[
 \epsilon_j\geq
   \frac{j-8}{7}-\frac{j-1}{6}
          +\frac{7j}{48}-\frac16
      =\frac{41j-384}{336}>0.
\]
At $j=8$, $v_7(8!)=1$ and \eqref{C22-unit} gives the
claimed equality.  Vanishing logarithmic derivatives simply
contribute no Bell monomials.

Write $F(z)=\sum_n a_nz^n$.  Its coefficient formula is
\[
 a_n=
 \sum_{\substack{m_j\geq0,\ \sum jm_j=n\\j\geq2\ {\rm even}}}
       \prod_j\frac{l_j^{m_j}}{m_j!(j!)^{m_j}}.
\]
For any nonzero summand, add $ns$ to its valuation and use
Legendre's formula.  The result is exactly
\begin{equation}\label{C22-partition}
 \sum_j m_j\epsilon_j
       +\frac16\sum_j s_7(m_j).
\end{equation}
Take $n=8\cdot7^r$.
The partition with $m_8=7^r$ and all other multiplicities
zero has value $1/6$ in \eqref{C22-partition}.
It is the only such partition: any other partition uses an
index $j\ne8$, so the first sum is strictly positive by
\eqref{C22-strict}; the second sum is at least $1/6$,
since some multiplicity is positive.
Thus the coefficient sum has a unique term of least valuation,
and the nonarchimedean triangle rule gives
\[
 v_7(a_n)=-\frac{7n}{48}+\frac16.
\]
Since the base-7 digits of $8\cdot7^r$ sum to two,
$v_7(n!)=(n-2)/6$.  Multiplying the coefficient by $n!$
therefore gives
\[
 v_7(F^{(n)}(0))=
    \frac{n-2}{6}-\frac{7n}{48}+\frac16
    =\frac{n}{48}-\frac16.
\]
The identity $D=d/dz$ in the additive coordinate proves
\eqref{C22-exact}.  Finally, if $c>1/48$, subtracting the
right side of \eqref{C22-exact} from $c n_r$ tends to
infinity.  This proves the asserted failure of every bounded-loss
estimate with the larger slope. \hfill$\square$

\paragraph{A second instance over $\Q(\sqrt{-3})$.}
The same phenomenon occurs on
$E:y^2=4x^3-4$ at $p=5$, with the nonzero point of
order two $T=(1,0)$ and $\omega=dx/y$.
This curve has good reduction at 5, CM by
$\Z[\zeta_3]$, and inert reduction of the CM prime.
The $\zeta_3$-symmetry forces $e_2^*=0$.
The same half-period argument makes the unit-normalized
actual Fock translate $F$ even, and again
$D^2\log F=-x(T+P(t))$ is integral.
Here $Dx=y$, $Dy=6x^2$ and
$D^4x=120x^3-48$, so
$l_6=-D^4x(T)=-72$ is a 5-adic unit.

For completeness the strict surplus calculation is also
explicit in this case.  Put
$\epsilon_j=v_5(l_j)-v_5(j!)+5j/24-1/4$.
The bound $v_5(l_j)\geq\lfloor(j-2)/5\rfloor$
gives $\epsilon_j>0$ for $j=2,4$;
$\epsilon_6=0$; and for all even $j\geq8$ it gives
\[
 \epsilon_j\geq
    \frac{j-6}{5}-\frac{j-1}{4}
          +\frac{5j}{24}-\frac14
       =\frac{19j-144}{120}>0.
\]
The identical Bell-partition argument, now with the
unique multiplicity $m_6=5^r$, yields
\[
 v_5\bigl(D^{6\cdot5^r}f_T(0)\bigr)
       =\frac{5^r-1}{4}
       =\frac{6\cdot5^r}{24}-\frac14.
\]
Thus the obstruction is present for this second CM field
and its stated actual theta normalization as well.

\paragraph{Positive tensor powers.}
In either example the same exact valuation formula holds
with $F$ replaced by $F^d$ for every positive integer $d$
prime to the displayed prime.  Indeed, the normalized
tensor section is still even and has constant term one,
and every logarithmic derivative $l_j$ is replaced by
$d l_j$.  Its valuation, and hence every surplus
$\epsilon_j$, is unchanged.  The same unique least
Bell monomial proves the same equality at every stated
index.  The obstruction therefore persists for these
actual positive tensor powers, with no adjustable
nonintegral scalar normalization.

\paragraph{What this decides.}
C19's weaker estimate for actual reduced theta tensors has
linear coefficient $1/(p^2-1)$, and this example shows that
coefficient is optimal for the class under consideration.
Equal-degree ratios still have the stronger intrinsic bound
because their common canonical frame cancels.
The example rules out transferring that stronger bound to
all individual sections by an unspecified torsion trace or
frame choice.  It makes no assertion that a differently
normalized archive theorem has these same degree, period,
or jet hypotheses; those data must be matched explicitly.
